\section{Mockup (High-Fidelity)}

Setelah wireframe divalidasi (Bab 3), langkah selanjutnya adalah mengembangkannya menjadi \textbf{mockup}---representasi visual \textit{high-fidelity} yang menunjukkan tampilan akhir halaman secara detail. Mockup mencakup warna, tipografi, gambar, ikon, dan spacing. Mockup memberikan gambaran yang sangat mirip dengan produk akhir, tetapi tanpa interaksi (tidak bisa diklik). Mockup dibuat setelah wireframe disetujui dan biasanya menggunakan alat desain seperti Figma, Adobe XD, atau Sketch \cite{figmaDocs}.

Mockup penting karena: (1) menjadi acuan pixel-perfect bagi pengembang saat implementasi; (2) membantu stakeholder memahami tampilan akhir sebelum development dimulai; dan (3) memastikan konsistensi visual di seluruh halaman.

\begin{table}[htbp]
\centering
\begin{tabular}{|P{3cm}|P{4.5cm}|P{4.5cm}|}
\hline
\textbf{Aspek} & \textbf{Wireframe (Low-Fi)} & \textbf{Mockup (High-Fi)} \\
\hline
Detail visual & Minimal (abu-abu, placeholder) & Lengkap (warna, font, gambar) \\
\hline
Tujuan & Struktur dan tata letak & Tampilan akhir yang akurat \\
\hline
Waktu pembuatan & Cepat (menit--jam) & Lebih lama (jam--hari) \\
\hline
Alat & Kertas, Balsamiq, Whimsical & Figma, Adobe XD, Sketch \\
\hline
Kapan dibuat & Tahap awal desain & Setelah wireframe disetujui \\
\hline
Interaktivitas & Tidak ada & Tidak ada (statis) \\
\hline
Revisi & Sangat mudah dan murah & Lebih memakan waktu \\
\hline
\end{tabular}
\caption{Perbandingan wireframe (low-fidelity) dan mockup (high-fidelity)}
\end{table}

\subsection{Dari Wireframe ke Mockup}

Pada tahap pengembangan wireframe menjadi mockup, desainer menambahkan elemen visual: palet warna yang sesuai, tipografi yang terbaca, ikon yang konsisten, gambar atau ilustrasi sesungguhnya, dan spacing yang presisi. Mockup harus mengikuti \textbf{design system} yang telah ditetapkan (jika ada) agar konsisten dengan halaman lain \cite{figmaDocs}.

\begin{contoh}
\textbf{Studi Kasus: Wireframe ke Mockup untuk Aplikasi Catatan}

\textbf{Wireframe:} Kotak abu-abu menunjukkan posisi header, daftar catatan (kartu), dan tombol ``Tambah''. Placeholder teks ``Lorem ipsum'' digunakan.

\textbf{Mockup:} Header menjadi biru tua (\code{\#1A237E}) dengan logo putih. Kartu catatan memiliki shadow halus, font Inter 14px untuk judul, warna teks \code{\#333}. Tombol FAB (Floating Action Button) berwarna hijau (\code{\#4CAF50}) dengan ikon ``+'' putih. Spacing antar kartu 12px.

\textbf{Pelajaran:} Wireframe memvalidasi \textit{struktur}, mockup memvalidasi \textit{tampilan}. Keduanya harus disetujui sebelum coding dimulai.
\end{contoh}

\begin{catatan}
Jangan langsung membuat mockup tanpa melalui wireframe terlebih dahulu. Memulai dari wire\-frame memungkinkan eksplorasi berbagai tata letak dengan cepat tanpa investasi waktu pada detail visual. Banyak pemula melakukan kesalahan ini---langsung mendesain mockup yang cantik, lalu harus mengubah seluruh strukturnya karena alur pengguna tidak tepat.
\end{catatan}
